% Consolidated Section 6: Conclusion and Future Scope.
% Standalone manuscript source. Replace sections/10-conclusion.tex with this file.
% The title, equation labels, results tables, and earlier sections are not modified.
% Original conclusion and accepted Sections 3--5 provide the scientific basis.

\section{Conclusion and Future Scope}
\label{sec:10}
\label{sec:conclusion}

The present study investigated how identity-aware biometric authorization can
be connected to encrypted ballot processing without transmitting a voter's
candidate selection through the inspected authorization request. Building
on \emph{Blockchain\textendash{}integrated IoT biometric architecture for
secure and decentralized electronic voting}
\citep{rifat2026biometric}, the proposed framework retains the inherited
fingerprint and facial authentication services while changing the voting
operations that follow authorization. A protected credential and an
election-specific membership witness are supplied to the voter-facing client,
where a local privacy worker constructs the encrypted ballot and its validity
proof. The architecture deliberately separates the existing Semaphore
registration and anchoring pathway from the custom Poseidon commitment tree
used for private-ballot membership. This implementation addresses the first
research question by demonstrating candidate-blind authorization at the
inspected service boundary alongside cryptographically checked ballot
submission. It does not establish that the credential custodian is unable to
identify the authorized voter or associate that voter with a public nullifier.

The encrypted voting relation combines membership in a frozen eligibility
set, an election-scoped nullifier, election-context binding, a valid one-hot
candidate selection, and consistency of the encrypted ballot coordinates.
Its fixed-depth Groth16 circuit uses a depth-17 Merkle path and 16 BabyJub
exponential-ElGamal candidate slots, with 73 public signals verified by the
voting contract. For an accepted submission, the contract rejects a reused
nullifier and adds the ciphertexts to an encrypted candidate-wise aggregate
without recording the selected candidate in plaintext. After the election
ends, the implemented three-of-five trustee procedure supports recovery of
aggregate counts through externally checked distributed-key-generation and
partial-decryption evidence. The contract checks the final result's state
and arithmetic conditions and records a commitment to the finalization
transcript. This establishes a concrete separation between on-chain
verification of encrypted-ballot validity and independent checking of tally
recovery. The recorded transcript hash alone does not prove correct
threshold decryption, and an independent auditor must obtain and verify the
complete evidence against the published encrypted aggregate.

The experimental evaluation characterizes the costs of this implemented
configuration rather than a general-purpose election deployment. Across
2,000 accepted encrypted-ballot transactions from 25 elections, mean gas
consumption was 1,314,307.123 gas and mean local submission-to-receipt
latency was 760.640~ms. A separate 20-trial fresh-chain campaign recorded
1,866,414 gas for the first accepted ballot, reflecting a different
execution state and measurement population. Candidate-count experiments
with 2, 4, 8, and 16 active candidates showed only a small gas difference
while retaining the same 16-slot ballot representation. With synthetic
eligible sets of 10 to 10,000 commitments and Merkle depth fixed at 17,
preparation costs increased as the populated set grew, whereas measured
mean Groth16 proving times remained between 5.292 and 5.423 seconds. These
findings address the measurement component of the second research question
within the controlled workloads. They do not measure complete
biometric-to-receipt sessions, public-network throughput, the cost of wider
circuits, or the performance of deeper Merkle relations. The historical
plaintext-voting gas figures are not a matched experimental baseline for
the more extensive encrypted-ballot operation.

The interpretation of these results depends on the threat boundaries
established for the prototype. The credential service can recover stored
voting secrets and derive identity-associated election nullifiers, even
though it does not receive the candidate choice through the inspected
authorization pathway. A compromised voting client may expose or alter the
selection before encryption, and a valid Groth16 proof cannot establish
that the interface displayed the voter's intended choice. Ballot secrecy
also depends on the integrity of encryption and protection of trustee
shares, since a sufficient coalition could target individually recorded
ciphertexts. Development-stage proving parameters, locally managed trustees,
and execution on Anvil limit the strength of deployment conclusions.
Functional checks and the measured performance campaigns do not establish
coercion resistance, receipt-freeness, independent trustee governance, or
production-election security assurance.

Future research should first strengthen the operational and cryptographic
boundaries that are most consequential for deployment. A credential design
with reduced custodian knowledge, stronger protection against terminal
compromise, and independently administered trustee shares would address
identity linkage and concentration of decryption capability. A
production-oriented evaluation would require a publicly reviewable proof
setup, accessible canonical decryption transcripts, and independently
reproducible verification of both the encrypted aggregate and the recovered
result. Further experiments should measure a complete biometric-to-receipt
session, varied client hardware, concurrent submission on distributed
networks, and threshold recovery under trustee failures. Testing additional
slot capacities and Merkle depths would distinguish fixed-configuration
performance from genuine circuit scaling. These investigations are needed
before extrapolating the present findings to operational elections. Within
its measured scope, the study demonstrates an integrated pathway for
candidate-blind biometric authorization, proof-verified encrypted voting,
and externally auditable threshold tally recovery.
